\section{The resulting noise floor}
\label{sec:noise}

Seven repetitions fall inside the later serving configuration. They scored
31.0, 31.0, 39.1, 34.5, 37.9, 37.9 and 34.5\%. \emph{Every headline figure
below is computed on the six rounds that exclude round 6, for the reason given
after Table~\ref{tab:loo}; seven-round values are given alongside so the
difference is always visible.} Seven of 29 tasks did not return the
same verdict every time (eight over all seven rounds), and the controls behave
as they should: 0 of 10 flips among tasks all published runs fail, 1 of 9 among
tasks all published runs pass, and 6 of 10 among the disagreed-on group.

\paragraph{The headline is not the range.} A best-minus-worst spread is an
order statistic: it cannot grow when a repetition is removed, it grows on
average with the number of repetitions, and a seven-round spread and a
six-round spread are therefore not comparable quantities. Our own sensitivity
check forced this correction. Table~\ref{tab:loo} recomputes everything with
each round dropped in turn: the spread is unmoved by dropping any round except
the one holding the maximum, while the minimum detectable difference varies
across all seven variants within a narrow band. Throughout, the
\emph{detectable difference at 95\%} is $1.96\sqrt{2}\,\sigma$, with $\sigma$
the resampled single-run standard deviation at the reweighted 117-task scale;
interval widths beside it are empirical 2.5--97.5\% spans of the same
distribution, hence not exact multiples (Appendix~\ref{app:intervals}).

\begin{table}[t]
\centering\small
\caption{Leave-one-round-out. Every figure recomputed from scratch on what
remains. The range responds only to removing the maximum; the detectable
difference is stable at 8.02--9.01 points and the comparison against published
entries never changes.}
\label{tab:loo}
\begin{tabular}{lccccccc}
\toprule
round dropped & 4 & 5 & \textbf{6} & 7 & 8 & 9 & 10 \\
\midrule
best $-$ worst (pp) & 8.10 & 8.10 & \textbf{6.90} & 8.10 & 8.10 & 8.10 & 8.10 \\
detectable difference (pp) & 9.00 & 8.52 & \textbf{8.15} & 8.42 & 8.02 & 8.76 & 9.01 \\
adjacent gaps within noise & 16/17 & 16/17 & 16/17 & 16/17 & 16/17 & 16/17 & 16/17 \\
\bottomrule
\end{tabular}
\end{table}

\paragraph{Round 6, and why the conservative number is the headline.} Round 6
contains the only infrastructure incident in the run: the provider returned
HTTP 401 for about an hour, 68 of the 95 silent restarts in this regime fall in it (110 across all ten rounds), and
all six runs that produced no verdict are in it. It is also the highest-scoring
round of the regime, on a reduced denominator of 23 tasks. We therefore report
the figures \emph{without} it as the headline, with the seven-round values
alongside:\footnote{The two scales move in \emph{opposite} directions when round
6 is removed: at the 29-task scale $\sigma$ rises from 5.87 to 5.92\,pp, at the
reweighted 117-task scale it falls from 3.13 to 2.94\,pp. They are computed
differently, and Appendix~\ref{app:sensitivity} traces the mechanism; the
detectable difference derives from the second scale. We report both because
reporting only the one that suits the argument is precisely the failure this
paper is about.}

\begin{center}\small
\begin{tabular}{lcc}
\toprule
 & 6 rounds (round 6 removed) & 7 rounds \\
\midrule
detectable difference at 95\% & \textbf{8.15 pp} & 8.67 pp \\
95\% interval at suite scale, width & \textbf{11.11 pp} & 12.82 pp \\
tasks that flipped & 7 / 29 & 8 / 29 \\
\bottomrule
\end{tabular}
\end{center}

\textbf{Removing the incident moves the result against us, and we report it
anyway.} The threshold falls from 8.67 to 8.15 points, which is a stricter bar:
the number of published pairs separated by less than it drops from 39 to 35.
We report the lower figure because it is the one measured on data containing no
infrastructure incident, not because it is the more favourable one.

Two things are worth stating precisely rather than glossing. First, the single
task that stops flipping, \texttt{DownloadSendReceiptTask}, does so at 1/7
$\rightarrow$ 0/6: it lost its one passing round, not its instability. Only 4 of
29 tasks change pass rate at all, and no task flips \emph{only} when round 6 is
removed. Second, imputing the six lost runs at those tasks' own rates in the
other rounds puts round 6 at 40.2\% rather than the 39.13\% recorded --- the
outage cost that round height rather than lending it any, because the tasks it
happened to cost pass more often than average (44.4\% against 34.5\%).
Appendix~\ref{app:sensitivity} has the full comparison.



\paragraph{Carrying it to the full suite.} The 29 tasks are deliberately
skewed, so their level means nothing on its own. Reweighting the three groups
to their true sizes in the 117-task GUI-only suite (98 disagreed-on, 9
always-passed, 10 always-failed) gives a width that is comparable to a
leaderboard number. Stated conditionally, because it is:
\emph{if the 88 disagreed-on tasks we did not measure behave like the 10 we
did}, a single run of the full suite falls within a 95\% interval about
11.1 points wide.

\paragraph{What this means for reading a leaderboard.} Under the resampling model above, and for this agent on this endpoint, two
single runs need to differ by 8.15 points before the gap exceeds what
repetition alone produces; to detect a real difference reliably rather than merely to notice it,
the gap must be larger still. Of the 17 adjacent pairs among the 18 published GUI-only entries, 16 are
separated by less than that, as are 35 of the 153 pairs overall (39 at the
seven-round threshold of 8.67). We state this as an indication of scale and not as a re-scoring: the
threshold was measured for one model on 29 tasks and is being applied to
entries produced by other agents and other models. The appropriate reading of a
gap below the threshold is that it \emph{warrants a second run before
citation}, not that the two systems are equal.

\paragraph{Are the flips scoring errors rather than variance?} It is the
obvious alternative explanation, and we checked it: we inspected the five
most suspicious of 111 ranked runs against the task definitions and found no
verifier defect --- every one was a genuine agent failure. Appendix~\ref{app:checker}
reports that negative result and why we had every incentive to find the
opposite.
